\chapter{Introducción}

El cambio climático, principalmente resultado de las actividades humanas, se ha convertido en uno de los retos más urgentes de nuestra era \parencite{allanSummaryPolicymakers2023}. Desde la época preindustrial, las concentraciones de gases de efecto invernadero (GEI) como el dióxido de carbono (CO\textsubscript{2}), el metano (CH\textsubscript{4}) y el óxido nitroso (N\textsubscript{2}O) han alcanzado niveles sin precedentes. En 2019, estas concentraciones llegaron a 410 partes por millón (ppm) para el CO\textsubscript{2}, 1866 partes por billón (ppb) para el CH\textsubscript{4} y 332 ppb para el N\textsubscript{2}O, según mediciones globales \parencite{ak-bhdWMOGreenhouseGas2021}. Este aumento ha intensificado el calentamiento global, provocando eventos climáticos extremos, alteraciones significativas en los ecosistemas y una rápida pérdida de biodiversidad.

Dentro de las fuentes antropogénicas de GEI, la gestión inadecuada de residuos orgánicos se destaca como un factor crítico, ya que su descomposición anaeróbica genera grandes cantidades de CH\textsubscript{4}, un gas con un potencial de calentamiento global 28 veces superior al del CO\textsubscript{2} en un horizonte de 100 años. Aunque los océanos y ecosistemas terrestres han actuado históricamente como sumideros naturales, absorbiendo aproximadamente el 56\,\% del CO\textsubscript{2} generado por el ser humano en las últimas décadas, esta capacidad muestra signos de saturación. Por tanto, es urgente desarrollar soluciones innovadoras y sostenibles para mitigar las emisiones actuales, particularmente aquellas vinculadas a la gestión de residuos orgánicos \parencite{ginerFoodLossWaste2025a}.

Las estrategias convencionales para la gestión de residuos orgánicos —como los rellenos sanitarios, la digestión anaerobia y el compostaje— enfrentan limitaciones significativas. Los rellenos sanitarios generan emisiones prolongadas y difícilmente controlables de metano durante años e incluso décadas; la digestión anaerobia, aunque permite capturar biogás, demanda infraestructura compleja y costosa; y el compostaje, a pesar de su mayor accesibilidad, conlleva emisiones de N\textsubscript{2}O y requiere períodos prolongados para lograr la estabilización del material \parencite{kongEvaluatingGreenhouseGas2012,nordahlGreenhouseGasAir2023}. En contextos agroindustriales de países en desarrollo —donde se concentra la mayor parte de los residuos orgánicos sin tratar—, estas alternativas suelen resultar inviables desde el punto de vista técnico o económico. Justamente esta brecha abre una ventana de oportunidad para tecnologías alternativas basadas en procesos biológicos que puedan ser controlados, escalables y adaptados a estas realidades.

Entre las alternativas emergentes, la bioconversión de residuos orgánicos mediante larvas de \textit{Hermetia illucens} (conocida como mosca soldado negra o BSF, por sus siglas en inglés) ha demostrado ser una tecnología prometedora y transformadora. A diferencia de los métodos convencionales, la BSF opera bajo un enfoque de economía circular: transforma residuos orgánicos en biomasa larval rica en proteínas y lípidos (utilizables como alimento animal o fuente de biocombustibles) y en fertilizante orgánico de alta calidad (\textit{frass}), todo mediante procesos biológicos que pueden ser diseñados y controlados \parencite{bermudezComprehensiveUtilizationBlack2023,Surendra2016}. Comparada con métodos tradicionales, la bioconversión con BSF presenta ventajas claras: requiere menor infraestructura, opera en tiempos más cortos (3--4 semanas frente a meses), funciona sin generación de lixiviados contaminantes y genera productos valiosos en lugar de residuos. Estudios preliminares sugieren que puede reducir significativamente las emisiones de N\textsubscript{2}O y CH\textsubscript{4} respecto al compostaje convencional \parencite{jenkinsProcessingPoultryManure2025}. Sin embargo, este potencial contrasta con un conocimiento aún limitado: la mayoría de las implementaciones se llevan a cabo sin monitoreo ambiental riguroso y sin una cuantificación fiable de los beneficios ambientales, económicos o de conversión que el sistema puede generar.

A pesar de los avances en la tecnología de bioconversión con BSF, persiste una brecha crítica en el conocimiento: la falta de marcos metodológicos robustos que permitan cuantificar y predecir, en tiempo real y a distintas escalas operativas, las emisiones de gases de efecto invernadero —en particular CO\textsubscript{2} y CH\textsubscript{4}— asociadas a estos procesos. Esta limitación constituye un obstáculo fundamental, ya que, sin datos fiables, reproducibles y contextualizados, la tecnología BSF no puede validarse como una solución climática verificable ni integrarse de forma creíble en estrategias de mitigación ambiental.

Esta tesis no tiene como objetivo diseñar un sistema de gestión ambiental, sino desarrollar una metodología transdisciplinaria que combine sensores de bajo costo, modelos dinámicos basados en ecuaciones diferenciales y algoritmos de aprendizaje automático para estimar en tiempo real las emisiones de CO\textsubscript{2} y CH\textsubscript{4} en procesos de bioconversión con BSF. La propuesta busca generar un desarrollo técnico accesible que, en el futuro, pueda ser adaptado por actores locales del Valle del Cauca —una región con intensa actividad agroindustrial y una creciente adopción de tecnologías de bioconversión— para tomar decisiones informadas, sin depender de infraestructura costosa o servicios externalizados. Al centrarse en la generación y validación local del conocimiento, esta aproximación responde a una necesidad práctica: cerrar la brecha entre la promesa tecnológica y su implementación justa, escalable y contextualmente pertinente.

Para cerrar esta brecha de conocimiento, se propone desarrollar un enfoque que vaya más allá de mediciones puntuales. Se busca \textbf{integrar tres componentes tecnológicos}: (1) \textbf{sensores con internet de las cosas (IoT) de bajo costo} que permitan monitoreo continuo en tiempo real de emisiones; (2) \textbf{modelos matemáticos basados en ecuaciones diferenciales ordinarias (EDO)} que capturen la dinámica del proceso en función de variables biológicas y ambientales (por ejemplo, biomasa larval, tasa de consumo de sustrato y respiración microbiana); y (3) \textbf{algoritmos de aprendizaje automático (AAA)} que permitan calibración dinámica y predicción adaptativa a escala horaria y diaria durante ciclos de bioconversión de 3--4 semanas.

Esta integración es crítica porque cada componente tiene limitaciones cuando opera aisladamente. Por ejemplo, \textbf{sensores NDIR pueden medir concentraciones de CO\textsubscript{2} con precisión}, pero sin un modelo matemático subyacente no es posible distinguir si los cambios observados en las mediciones corresponden a comportamientos normales y predecibles del sistema o a anomalías que requieren intervención. \textbf{Modelos EDO pueden predecir el comportamiento teórico del sistema} basándose en principios de crecimiento microbiano y balance de energía, pero sin datos en tiempo real no se calibran adecuadamente y sus predicciones pueden diferir de la realidad operativa. \textbf{Algoritmos de aprendizaje automático aprenden patrones complejos en datos históricos}, pero sin una estructura matemática mecanística subyacente pueden hacer predicciones sin sentido físico o generalizar incorrectamente a nuevas condiciones. Solo la combinación de los tres componentes permite un sistema que sea simultáneamente preciso (datos en tiempo real), interpretable (modelo matemático) y adaptativo (aprendizaje automático).

El objetivo de esta tesis es \textbf{desarrollar y validar un sistema  híbrido que integre sensores con internet de las cosas (IoT), modelos matemáticos y algoritmos de aprendizaje automático (AAA)} para predecir, en tiempo real y con la mayor precisión posible, las emisiones de CO\textsubscript{2} y CH\textsubscript{4} durante la bioconversión de residuos orgánicos con BSF. Este sistema busca medir emisiones, comprender su dinámica y predecir su comportamiento.

No obstante, la propuesta enfrenta \textbf{limitaciones técnicas reales}. El funcionamiento confiable de sensores en ambientes de alta humedad presenta desafíos significativos de calibración y estabilidad; la disponibilidad inicial limitada de datos podría restringir la capacidad de los modelos de aprendizaje automático; y la escalabilidad del sistema en contextos rurales con infraestructura limitada (conectividad inestable, acceso limitado a energía, capacitación técnica insuficiente) requiere soluciones de conectividad y operación robustas. Además, debe considerarse críticamente el uso de tecnologías emergentes como la inteligencia artificial en estos procesos, evaluando tanto sus oportunidades como sus riesgos en términos de dependencia tecnológica de usuarios rurales, posibles sesgos en las predicciones basadas en datos limitados y requisitos energéticos que pueden ser problemáticos en contextos rurales aislados.

A pesar de estas limitaciones, la propuesta se sitúa en una intersección estratégica entre biotecnología, electrónica, modelado matemático e inteligencia artificial. Este enfoque interdisciplinario responde a una necesidad urgente: transformar la bioconversión con BSF de una solución productiva aislada hacia un sistema inteligente, trazable y ambientalmente responsable que contribuya efectivamente a la mitigación del cambio climático y a la economía circular, con beneficios equitativos para pequeños y medianos productores.

La importancia de este trabajo radica, entonces, en su doble apuesta: por un lado, contribuir al rigor científico mediante la integración de modelado dinámico y datos empíricos en condiciones reales; por otro, ejercer una crítica constructiva a la lógica tecno-solucionista que suele acompañar a las propuestas de inteligencia artificial en el sur global. Al reconocer que los modelos de IA reflejan los contextos en los que se entrenan —y que rara vez incluyen la variabilidad biológica, climática o socioeconómica de sistemas agroindustriales locales—, esta investigación se compromete con un enfoque epistemológicamente responsable. No se trata solo de medir gases, sino de construir herramientas que respeten, integren y potencien los saberes existentes, evitando que la digitalización de la sostenibilidad se convierta en un nuevo vector de exclusión o dependencia tecnológica.

\textbf{Nota sobre el alcance y responsabilidad:} 
Si bien esta investigación se centra en los componentes técnicos 
del sistema, reconoce explícitamente que la implementación de 
tecnologías emergentes en contextos rurales exige consideraciones 
más amplias sobre acceso a datos, pertinencia cultural y gobernanza 
local que quedan fuera del presente alcance pero son críticas para 
una adopción verdaderamente sostenible.